% title:          Sum of two nilpotent elements
% area:           ring-theory
% methods:        construction
% difficulty:
% difficulty-ai:  easy
% status:
% review:         none
% --- history: all optional, leave EMPTY if unknown ---
% origin:         folklore
% origin-date:
% origin-number:
% also-in:        book, jarnik
% taken-from:     VJIMC 2006 official problems and solutions (Category II, Problem 1)
% references:     Folklore: in a commutative ring (identity not needed) the nilpotent elements form an ideal, the nilradical (binomial theorem); without commutativity a sum of two nilpotent elements need not be nilpotent (2x2 matrices). Standard textbook statement: M. F. Atiyah, I. G. Macdonald, Introduction to Commutative Algebra, Addison-Wesley 1969, Chapter 1, Proposition 1.7, p. 5 (proof via (x+y)^(m+n-1) = 0; 'ring' there means commutative with identity, p. 1) - https://u.cs.biu.ac.il/~plotkin/resources/MFAtiyah_IGMacDonald_IntroToCommAlgebra.pdf ; earliest explicit source found (it calls the fact 'klar', so probably not the first): E. Noether, Idealtheorie in Ringbereichen, Math. Ann. 83 (1921), 24-66, §4, p. 38 (commutative ring, identity not assumed: for a primary ideal Q the elements with a power in Q form an ideal) and §9, p. 56, footnote 38 (in non-commutative rings p1^λ1, p2^λ2 in Q do not give (p1 - p2)^(λ1+λ2) in Q), doi:10.1007/BF01464225 - https://gdz.sub.uni-goettingen.de/download/pdf/PPN235181684_0083/LOG_0008.pdf ; these two were opened ; seen only as Google Books snippets, not opened: N. H. McCoy, Rings and Ideals, 1948 ('The set of all nilpotent elements is merely the radical of the zero ideal, and hence is an ideal'); B. L. van der Waerden, Algebra II, 4th ed., Springer 1959, 'Allgemeine Idealtheorie der kommutativen Ringe' ('In einem kommutativen Ring A bilden die nilpotenten Elemente ein Ideal N'); J. Barshay, Topics in Ring Theory, 1969, Exercise 8-4 ('Show by example that a sum of nilpotent elements need not, however, be nilpotent') ; competition: 16th Vojtěch Jarník International Mathematical Competition (VJIMC), Ostrava, 29 March 2006, Category II, Problem 1 - https://vjimc.osu.cz/storage/uploads/j16problems2.pdf ; solution 'by Stijn Cambie' in the VJIMC solutions file, 5th of 8 pages (PDF made 12 Aug 2014, so a later write-up, not a 2006 jury text; word for word the club copy, including the formula (u+v)^n = [[0,(-2)^n],[2^n,0]], true only for n = 1 mod 4) - https://vjimc.osu.cz/storage/uploads/j16solutions.pdf
% history:        ai
%
% AI-WRITTEN, NEEDS HUMAN REVIEW (2026-09-29): both the statement and the solution were written by AI (Claude).
% The statement makes the official VJIMC 2006 text rigorous (ring axioms, powers, commutative, no identity assumed).
% The solution follows the solution of the club copy (by Stijn Cambie) with these changes: in (a), the binomial formula,
% which uses u^0 (undefined without an identity), is replaced by the expansion of the product into words (indexed by
% {0,1}^N, so that u = v is covered too), and the facts the copy only asserts (y0 = 0, u^t = 0 for t >= n) are proved;
% in (b), the copy claims (u+v)^n = [[0,(-2)^n],[2^n,0]], which holds only for n = 1 (mod 4): u + v = 2J with J^2 = -I,
% e.g. (u+v)^2 = -4I; the solution uses (u+v)^2 = -4I, and checks that uv and vu differ.
% Restyled on 2026-10-01 after problem-bank/writing-style.md (the style of the fully solved problems); the mathematics is unchanged.
\begin{problem}
In this problem, a ring \( \left(R, +, \cdot, 0_R\right) \) is not necessarily unital: \( \left(R, +\right) \) is an abelian group with neutral element \( 0_R \), and the multiplication \( \cdot \), written \( xy := x \cdot y \), is associative and distributes over \( + \) on both sides. The ring is commutative if \( xy = yx \) for all \( x, y \in R \). For \( x \in R \) and \( k \in \mathbb{N}_{>0} \), define the power \( x^{k} \) by \( x^{1} := x \) and \( x^{k+1} := x^{k} x \); an element \( x \in R \) is nilpotent if \( x^{k} = 0_R \) for some \( k \in \mathbb{N}_{>0} \).

(a) Let \( R \) be a commutative ring (with or without identity) and let \( u, v \in R \) be nilpotent. Prove that \( u + v \) is nilpotent.

(b) Give an example of a ring \( R \) (by (a), necessarily non-commutative) and of nilpotent elements \( u, v \in R \) such that \( u + v \) is not nilpotent.
\end{problem}
\begin{solution}
We begin by collecting some facts that are valid in any ring \( R \), commutative or not. Let \( x \in R \). Since \( \cdot \) is associative, we have \( x^{k} x^{l} = x^{k+l} \) for all \( k, l \in \mathbb{N}_{>0} \), by induction on \( l \): for \( l = 1 \), \( x^{k} x^{1} = x^{k} x = x^{k+1} \) by definition, and if the formula holds for some \( l \geq 1 \), then
\[
x^{k} x^{l+1} = x^{k} \left(x^{l} x\right) = \left(x^{k} x^{l}\right) x = x^{k+l} x = x^{k+l+1}.
\]
Hence, by induction on \( k \), we have \( \left(x^{2}\right)^{k} = x^{2k} \) for every \( k \in \mathbb{N}_{>0} \): for \( k = 1 \), \( \left(x^{2}\right)^{1} = x^{2} \), and if the formula holds for some \( k \geq 1 \), then
\[
\left(x^{2}\right)^{k+1} = \left(x^{2}\right)^{k} x^{2} = x^{2k} x^{2} = x^{2k+2}.
\]
Moreover, \( y \cdot 0_R = 0_R = 0_R \cdot y \) for every \( y \in R \): since \( \cdot \) distributes over \( + \), we have \( y \cdot 0_R = y \left(0_R + 0_R\right) = y \cdot 0_R + y \cdot 0_R \), and adding \( -\left(y \cdot 0_R\right) \) to both sides gives \( y \cdot 0_R = 0_R \); symmetrically, \( 0_R \cdot y = 0_R \). Thus, if \( x^{n} = 0_R \) for some \( n \in \mathbb{N}_{>0} \), then \( x^{t} = 0_R \) for every \( t \geq n \), since for \( t > n \) we have \( x^{t} = x^{t-n} x^{n} = x^{t-n} \cdot 0_R = 0_R \).
\begin{enumerate}
    \item[(a)] Let \( n, m \in \mathbb{N}_{>0} \) be such that \( u^{n} = 0_R \) and \( v^{m} = 0_R \), and define \( N := n + m \), \( w_{0} := u \) and \( w_{1} := v \). Since \( \cdot \) distributes over \( + \), expanding the product of the \( N \) factors \( u + v = w_{0} + w_{1} \) gives
    \[
    \left(u + v\right)^{N} = \sum_{s \in \left\{0, 1\right\}^{N}} \prod_{i \in N} w_{s_{i}},
    \]
    where each product is taken in the order of the natural number \( N = \left\{ i \in \mathbb{N} \mid i < N \right\} \) (one term for each choice of a summand in each factor; this also holds when \( u = v \)). Fix \( s \in \left\{0, 1\right\}^{N} \) and let \( r \) be the number of indices \( i \in N \) with \( s_{i} = 0 \). Since \( R \) is commutative, the factors can be reordered, so the product equals \( u^{r} v^{N-r} \) if \( 1 \leq r \leq N - 1 \), \( u^{N} \) if \( r = N \), and \( v^{N} \) if \( r = 0 \) (no identity element is needed). We show that it is \( 0_R \) in every case.
    \\\\
    If \( r \geq n \), then \( u^{r} = 0_R \) by the above, so the product is \( u^{N} = 0_R \) when \( r = N \), and \( u^{r} v^{N-r} = 0_R \cdot v^{N-r} = 0_R \) otherwise.
    \\\\
    If \( r \leq n - 1 \), then \( N - r \geq m + 1 \), so \( v^{N-r} = 0_R \) by the above, and the product is \( v^{N} = 0_R \) when \( r = 0 \), and \( u^{r} v^{N-r} = u^{r} \cdot 0_R = 0_R \) otherwise.
    \\\\
    Hence, all \( 2^{N} \) products vanish, so \( \left(u + v\right)^{N} = 0_R \) and \( u + v \) is nilpotent.

    \item[(b)] Take the ring \( \left(\mathbb{Z}^{2 \times 2}, +, \cdot, \mathbf{0}_{2 \times 2}\right) \) of \( 2 \times 2 \) integer matrices with the usual addition and multiplication (its identity is \( I_{2} \)), and
    \[
    u := \begin{pmatrix} 1 & -1 \\ 1 & -1 \end{pmatrix}, \quad v := \begin{pmatrix} -1 & -1 \\ 1 & 1 \end{pmatrix}.
    \]
    A direct computation gives
    \[
    u^{2} = \begin{pmatrix} 1 - 1 & -1 + 1 \\ 1 - 1 & -1 + 1 \end{pmatrix} = \mathbf{0}_{2 \times 2}, \quad v^{2} = \begin{pmatrix} 1 - 1 & 1 - 1 \\ -1 + 1 & -1 + 1 \end{pmatrix} = \mathbf{0}_{2 \times 2},
    \]
    so \( u \) and \( v \) are nilpotent, while
    \[
    u + v = \begin{pmatrix} 0 & -2 \\ 2 & 0 \end{pmatrix} \quad \text{and} \quad \left(u + v\right)^{2} = \begin{pmatrix} -4 & 0 \\ 0 & -4 \end{pmatrix} = -4 I_{2}.
    \]
    Suppose that \( \left(u + v\right)^{k} = \mathbf{0}_{2 \times 2} \) for some \( k \in \mathbb{N}_{>0} \). Then \( \left(u + v\right)^{2k} = \mathbf{0}_{2 \times 2} \) by the above (as \( 2k \geq k \)), whereas, again by the above,
    \[
    \left(u + v\right)^{2k} = \left(\left(u + v\right)^{2}\right)^{k} = \left(-4 I_{2}\right)^{k} = \left(-4\right)^{k} I_{2} \neq \mathbf{0}_{2 \times 2},
    \]
    which is a contradiction. Hence, \( u + v \) is not nilpotent, as required. By (a), this ring is necessarily non-commutative; indeed,
    \[
    uv = \begin{pmatrix} -2 & -2 \\ -2 & -2 \end{pmatrix} \neq \begin{pmatrix} -2 & 2 \\ 2 & -2 \end{pmatrix} = vu.
    \]
\end{enumerate}
\end{solution}
